\chapter{Il sistema proposto}
\label{ch:core}

% BOZZA — architettura del sistema, da raffinare man mano che il prototipo
% cresce. Titolo del capitolo provvisorio.

\section{Architettura generale}
% La pipeline in tre fasi:
% 1. Contesto: al modello si fornisce il file sorgente, oppure una
%    rappresentazione di più alto livello (riassunto AST), oppure il file
%    più uno scheletro di test generato da Klara.
% 2. Generazione iterativa: ogni test generato viene eseguito (pytest);
%    gli errori vengono restituiti al modello, che itera fino a ottenere
%    test funzionanti.
% 3. Feedback di coverage: si calcola la coverage (coverage.py) e si
%    passano al modello le righe non coperte, chiedendo test mirati.

\section{Estrazione del contesto strutturale}
% Dagli esperimenti con la libreria ast: per ogni funzione si estraggono
% nome e argomenti, docstring, numero di return e di if (legato al numero
% di cammini da coprire), eccezioni sollevate, righe di inizio/fine.
% Il riassunto è molto più compatto del file: utile per file grandi
% (limiti di contesto, costo per token). lineno/end_lineno permettono di
% mappare le righe non coperte alla funzione corrispondente (collega la
% fase 3 alla fase 1).

\section{Scheletri di test con Klara}
% Ruolo: fornire lo scheletro iniziale da cui l'LLM parte. Limiti:
% supporta un sottoinsieme ristretto di Python (crash su list
% comprehension, eccezioni, Python moderno; non installabile su
% Python >= 3.12). Decisione sul suo ruolo da prendere con la relatrice.

\section{Il ciclo di feedback}
% Il cuore del sistema: genera -> esegui -> misura -> feedback -> migliora.
% Stato del prototipo: script di orchestrazione (ciclo.py) semi-automatico,
% con passaggio manuale verso l'LLM in attesa dell'accesso API.

\section{Criterio di arresto}
% Dall'esperimento 5: il loop non può fermarsi "al 100%" perché esistono
% righe irraggiungibili (codice morto o difensivo). Criterio: plateau di
% coverage o numero massimo di iterazioni. Le righe che restano scoperte
% dopo il plateau hanno valore diagnostico (candidate a essere codice
% morto); la dimostrazione formale di irraggiungibilità è il mestiere dei
% solver SMT come Z3 (possibile punto di contatto con Klara).
